\documentclass[runningheads]{llncs}

\usepackage{graphicx}
\usepackage{amsmath}
\usepackage{amssymb}
\usepackage{booktabs}
\usepackage{url}

\title{ASTRID: Local-First, Behavior-Aware Dataset Reliability Assessment for Industrial AI}

\author{Jorge Martinez Gil}
\institute{Affiliation to be added}

\begin{document}
\maketitle

\begin{abstract}
Industrial AI systems increasingly fail not because a model class is unavailable,
but because training and evaluation data contain defects that are only partially
visible to conventional quality checks. Missing values, duplicate records,
cross-split leakage, temporal drift, personally identifiable information, and
subgroup imbalance can all pass through ordinary data preparation pipelines while
creating brittle or misleading machine learning workflows.

This paper introduces ASTRID, a local-first framework for behavior-aware dataset
reliability assessment. ASTRID treats dataset reliability as an auditable evidence
vector over quality, security, reliability, robustness, fairness, and transparency
dimensions, and converts this evidence into diagnostic findings, recommendations,
and policy-gate outcomes. The framework is implemented both as an interactive
tool and as a headless analyzer suitable for reproducible experiments and CI-style
dataset gates.

We evaluate the current tabular implementation through a controlled fault-injection
benchmark with 180 audits: five random seeds, six fault families, and six severity
levels. Clean controls produce zero primary detections. ASTRID detects split
leakage and PII at 1\% injected severity, duplicates at 5\%, fairness disparity at
10\%, and drift and missingness at 20\% under the configured policy thresholds.
The health score decreases under every injected fault family, and governance
signals--recommendations and policy failures--track the same threshold crossings.
The mean audit time is 0.229 seconds per run on 4,000-row synthetic tabular
datasets.

The results support a conservative but useful claim: behavior-aware pre-training
dataset audits can expose measurable reliability risks before they enter model
development workflows. The paper also identifies the remaining experiments needed
for a full top-tier empirical claim, including real public benchmarks, baselines,
downstream model-degradation studies, and multimodal fault injection.
\keywords{Dataset reliability \and Data-centric AI \and Industrial AI \and Data quality \and Responsible AI \and MLOps}
\end{abstract}

\section{Introduction}

Industrial AI systems are now embedded in predictive maintenance, manufacturing
inspection, process optimization, logistics, and operational decision support.
In these settings, the reliability of a deployed model is strongly constrained by
the reliability of the data used to train, validate, and govern it. Seemingly
ordinary dataset defects can have operational consequences: duplicated rows can
inflate validation scores, cross-split leakage can invalidate performance
estimates, temporal drift can hide seasonal failure modes, and subgroup imbalance
can make aggregate metrics appear stable while specific sites or operating regimes
are underserved.

Conventional dataset checks are necessary but insufficient. Database and data
quality practice provides measures such as completeness, consistency, validity,
timeliness, and uniqueness \cite{iso25012,schelter2018}. ML engineering practice
adds schema validation, distribution checks, and monitoring \cite{breck2019,
polyzotis2017}. These methods are valuable, but they often answer a narrow
question: does the dataset satisfy a set of static constraints? Industrial AI
requires an additional question: does the dataset provide enough evidence to
support reliable model development, evaluation, and governance?

We call this perspective \emph{behavior-aware dataset reliability}. The word
``behavior'' is used deliberately. The aim is not merely to describe a table or
file collection, but to expose data conditions that are known to alter the
behavior of ML workflows: leakage, drift, sensitive information exposure,
fairness disparity, missingness, and duplicate records. A behavior-aware audit
therefore combines descriptive measurement with diagnostic interpretation and
policy consequences.

This paper presents ASTRID, a local-first framework for dataset reliability
assessment in industrial AI. ASTRID is designed for environments where data often
cannot leave the organization, where audit artifacts must be reproducible, and
where dataset gates should be understandable by data engineers, ML practitioners,
and governance reviewers. The current implementation supports tabular reliability
audits through both an interactive Streamlit interface and a headless Python API.
The architecture is extensible to other modalities, but the empirical claims in
this paper are intentionally restricted to the tabular implementation that has
been tested.

The contributions are:

\begin{itemize}
    \item A behavior-aware reliability model that represents dataset risk as a
    multi-dimensional evidence vector rather than as a single static data quality
    statistic.
    \item The ASTRID framework, a local-first implementation that produces
    diagnostic findings, recommendations, reports, and configurable policy-gate
    decisions.
    \item A reproducible tabular fault-injection benchmark covering missingness,
    duplicates, cross-split leakage, temporal drift, PII exposure, and fairness
    disparity.
    \item An empirical evaluation over 180 audits showing threshold-consistent
    detection, monotone or near-monotone metric response, health-score degradation,
    recommendation coverage, policy-gate alignment, and low audit runtime.
\end{itemize}

\section{Related Work}

\subsection{Data Quality and Validation}

Data quality has been studied extensively in databases, information systems, and
software quality standards. Common dimensions include completeness, accuracy,
consistency, timeliness, and uniqueness \cite{iso25012}. Modern ML data validation
systems operationalize these ideas through schema constraints, statistical
summaries, and anomaly checks \cite{breck2019,schelter2018}. ASTRID builds on this
tradition but differs in emphasis: its output is not only a set of constraint
violations, but an audit record organized around ML-relevant reliability risks
and governance decisions.

\subsection{Data-Centric AI and Dataset Documentation}

Data-centric AI emphasizes that improvements to data can be as important as, and
often more durable than, changes to model architecture \cite{sambasivan2021}.
Dataset documentation practices such as datasheets and model cards improve
transparency by requiring explicit statements about collection, intended use,
limitations, and evaluation context \cite{gebru2021,mitchell2019}. ASTRID is
complementary: it provides executable evidence that can populate such
documentation and can be regenerated as datasets evolve.

\subsection{Robustness, Leakage, and Behavioral Testing}

Robustness research studies how systems respond to perturbations, distribution
shift, and adversarial or noisy conditions. Behavioral testing has shown that
targeted test suites can reveal failures that aggregate accuracy hides
\cite{ribeiro2020}. Data leakage research similarly demonstrates that models can
appear highly accurate when train-test separation is invalid \cite{kaufman2012}.
ASTRID applies a related testing philosophy before model training: the dataset is
audited for conditions likely to undermine trustworthy evaluation and deployment.

\subsection{Production ML Governance}

Production ML systems accumulate hidden technical debt through data dependencies,
configuration, monitoring gaps, and feedback loops \cite{sculley2015,
amershi2019}. Governance frameworks such as the NIST AI Risk Management Framework
emphasize measurement, documentation, and risk controls \cite{nist2023}. ASTRID
is positioned as a practical measurement layer for dataset-level controls in
MLOps pipelines.

\section{Behavior-Aware Dataset Reliability}

\subsection{Dataset Evidence Model}

Let $D$ denote a dataset with features $X$, optional labels $Y$, metadata $M$,
split assignments $S$, time information $T$, and optional group attributes $G$:
\[
D = (X, Y, M, S, T, G).
\]
ASTRID computes an evidence vector
\[
E(D) = \{E_q, E_s, E_r, E_b, E_f, E_t\},
\]
where the components correspond to data quality, security/privacy, reliability,
robustness, fairness, and transparency. Each component is itself a collection of
measurable indicators, such as missing-value rate, duplicate-row rate,
cross-split row-hash overlap, temporal drift statistics, PII flags, and subgroup
positive-rate disparity.

\subsection{Health Score and Policy Gate}

For reporting, ASTRID maps evidence to a bounded health score:
\[
H(D) = 100 - \sum_{j=1}^{m} w_j p_j(D),
\]
where $p_j(D)$ is the normalized penalty associated with finding $j$, and $w_j$
is its configured weight. The health score is a decision-support index, not a
proof of safety. It is paired with a policy gate
\[
G(D; \theta) \in \{\mathrm{pass}, \mathrm{fail}\},
\]
where $\theta$ contains operational thresholds, for example maximum duplicate
rate, maximum missingness, maximum drift statistic, allowed PII findings, and
maximum positive-rate disparity.

This separation is important. A continuous score supports ranking and trend
analysis, while the policy gate supports concrete release decisions. A dataset
can have a moderate health score but still fail because it contains a critical
finding such as leakage or PII.

\subsection{Fault-Oriented Reliability}

ASTRID evaluates reliability through fault-oriented diagnostics. Given a family
of defect operators $\mathcal{T} = \{T_1,\ldots,T_k\}$ and severities
$\alpha \in A$, a controlled benchmark produces datasets
\[
D_{i,\alpha} = T_i(D, \alpha).
\]
The expected behavior is that the primary metric associated with fault family
$i$ increases with $\alpha$, and that health score, recommendations, and policy
decisions respond when thresholds are crossed. This is weaker than proving
downstream model degradation, but it is a necessary first validation step for a
dataset reliability tool: the tool must detect known injected defects before it
can be trusted on unknown real defects.

\section{The ASTRID Framework}

\subsection{Architecture}

ASTRID consists of five layers:

\begin{enumerate}
    \item \textbf{Ingestion:} load user-provided datasets and identify optional
    semantic columns such as label, split, timestamp, group, and identifier.
    \item \textbf{Metric extraction:} compute modality-aware indicators for
    completeness, uniqueness, leakage, drift, PII, fairness, and transparency.
    \item \textbf{Risk interpretation:} map indicators to findings, severity
    levels, and recommendations.
    \item \textbf{Governance gate:} evaluate configurable thresholds and produce
    pass/fail decisions.
    \item \textbf{Reporting:} generate human-readable summaries and machine-readable
    JSON artifacts for audit history and reproducibility.
\end{enumerate}

The implementation is local-first: analysis runs on the user's machine and does
not require uploading datasets to an external service. This design is important
for industrial data, which often contains operational, confidential, or regulated
information.

\subsection{Current Scope}

The current tested implementation focuses on tabular datasets. It supports
numeric and categorical columns, optional text fields for PII scanning, explicit
train/validation/test split columns, timestamp columns for temporal drift, group
columns for fairness checks, and identifier columns for duplicate and leakage
diagnostics. The broader ASTRID interface is designed to accommodate time-series
and image analyzers, but the experiments in this paper do not claim empirical
validation for those modalities.

\subsection{Outputs}

ASTRID produces four output types. First, scalar metrics expose the measured
dataset conditions. Second, a health score summarizes the overall audit result.
Third, recommendations convert findings into concrete remediation actions. Fourth,
the policy gate determines whether the dataset satisfies configured reliability
requirements. These outputs target different users: engineers need metrics,
practitioners need recommendations, and governance stakeholders need pass/fail
evidence with a reproducible trail.

\section{Experimental Design}

\subsection{Research Questions}

The evaluation addresses three questions:

\begin{itemize}
    \item \textbf{RQ1:} Does ASTRID detect controlled tabular reliability faults
    when they cross configured thresholds?
    \item \textbf{RQ2:} Do ASTRID's primary metrics and health score respond
    coherently as injected severity increases?
    \item \textbf{RQ3:} Do recommendations and policy gates align with detected
    reliability risks, and is runtime practical for CI-style use?
\end{itemize}

\subsection{Synthetic Industrial Tabular Dataset}

The benchmark uses a synthetic industrial tabular dataset with 4,000 clean rows
per seed. Each row contains a sample identifier, timestamp, split assignment,
site identifier, numeric sensor features, an operator-note text field, and a
binary target. The columns are chosen to exercise common industrial AI metadata:
splits for leakage checks, timestamps for drift checks, sites for subgroup
analysis, identifiers for duplicate detection, and text for PII scanning.

Synthetic data is appropriate for this first validation because the ground truth
fault type and severity are exactly controlled. It is not a substitute for real
industrial benchmarks; that limitation is addressed in Section~\ref{sec:limits}.

\subsection{Fault Injection Protocol}

Each experiment starts from a clean dataset and injects one fault family at one
severity level. Faults are independent, so the benchmark measures isolated
response rather than interacting real-world failures. We use six fault families:

\begin{itemize}
    \item \textbf{Missingness:} random missing values in numeric sensor columns.
    \item \textbf{Duplicates:} exact duplicate row insertion.
    \item \textbf{Split leakage:} duplicated samples copied across train/test
    split boundaries.
    \item \textbf{Drift:} numeric feature shifts in the final temporal slice.
    \item \textbf{PII:} synthetic email-like strings inserted into operator notes.
    \item \textbf{Fairness:} group-dependent positive-rate disparity.
\end{itemize}

The severity grid is
\[
A=\{0, 0.01, 0.05, 0.10, 0.20, 0.40\}.
\]
The zero-severity setting is a clean control for each fault family and seed. The
experiment uses five seeds: 7, 11, 19, 23, and 31. This yields
$6 \times 6 \times 5 = 180$ audits.

\subsection{Detection Rules}

Table~\ref{tab:detection-rules} lists the primary metric and detection rule for
each fault family. These rules are aligned with the current ASTRID policy
configuration and are therefore best interpreted as threshold-response tests, not
as claims that the thresholds are globally optimal.

\begin{table}
\centering
\caption{Fault families, primary metrics, and detection rules.}
\label{tab:detection-rules}
\begin{tabular}{lll}
\toprule
Fault family & Primary metric & Detection rule \\
\midrule
Missingness & Overall missing rate & $> 0.05$ \\
Duplicates & Exact duplicate row rate & $> 0.01$ \\
Split leakage & Cross-split row-hash rate & $> 0$ \\
Drift & Max numeric KS statistic & $> 0.30$ \\
PII & PII flagged columns & $> 0$ \\
Fairness & Positive-rate disparity & $> 0.20$ \\
\bottomrule
\end{tabular}
\end{table}

\subsection{Reproducibility}

The benchmark is implemented as a headless experiment runner. The canonical run
was executed with the following configuration:
\[
\texttt{rows}=4000,\quad
\texttt{seeds}=7,11,19,23,31,\quad
\texttt{severities}=0,0.01,0.05,0.10,0.20,0.40.
\]
The run generated per-audit JSON reports, a summary CSV, an aggregate CSV, a
results JSON file, and SVG figures for score, primary metrics, detection,
recommendation coverage, and policy-gate behavior.

\section{Results}

\subsection{RQ1: Detection Behavior}

Table~\ref{tab:detection-rate} reports mean detection rate across the five seeds.
Clean controls have zero primary detections for every fault family. Split leakage
is detected at the smallest nonzero severity because any duplicated row across
split boundaries invalidates evaluation. PII is detected at 1\% severity for
three of five seeds and at all higher severities. Duplicate rows are consistently
detected from 5\% severity. Fairness disparity is detected for four of five seeds
at 10\% severity and all seeds at 20\%. Drift and missingness are detected from
20\% severity, consistent with their configured thresholds.

\begin{table}
\centering
\caption{Detection rate by fault family and severity, averaged across five seeds.}
\label{tab:detection-rate}
\scriptsize
\begin{tabular}{lrrrrrr}
\toprule
Fault family & 0 & 0.01 & 0.05 & 0.10 & 0.20 & 0.40 \\
\midrule
Drift & 0.00 & 0.00 & 0.00 & 0.00 & 1.00 & 1.00 \\
Duplicates & 0.00 & 0.00 & 1.00 & 1.00 & 1.00 & 1.00 \\
Fairness & 0.00 & 0.00 & 0.00 & 0.80 & 1.00 & 1.00 \\
Missingness & 0.00 & 0.00 & 0.00 & 0.00 & 1.00 & 1.00 \\
PII & 0.00 & 0.60 & 1.00 & 1.00 & 1.00 & 1.00 \\
Split leakage & 0.00 & 1.00 & 1.00 & 1.00 & 1.00 & 1.00 \\
\bottomrule
\end{tabular}
\end{table}

\subsection{RQ2: Health Score Response}

Table~\ref{tab:score-mean} reports mean ASTRID health score. The clean score is
91.4 across fault families because the zero-severity condition is the same clean
dataset protocol. The score decreases under every injected fault family. PII and
split leakage cause large early penalties because they represent critical
governance risks. Missingness and drift degrade more gradually until their
thresholds are crossed.

\begin{table}
\centering
\caption{Mean ASTRID health score by fault family and severity. Higher is better.}
\label{tab:score-mean}
\scriptsize
\begin{tabular}{lrrrrrr}
\toprule
Fault family & 0 & 0.01 & 0.05 & 0.10 & 0.20 & 0.40 \\
\midrule
Drift & 91.4 & 91.4 & 89.8 & 84.2 & 78.2 & 78.2 \\
Duplicates & 91.4 & 90.4 & 85.8 & 80.4 & 79.2 & 79.4 \\
Fairness & 91.4 & 91.6 & 89.8 & 87.6 & 83.4 & 81.8 \\
Missingness & 91.4 & 91.2 & 90.4 & 89.2 & 87.0 & 83.2 \\
PII & 91.4 & 76.4 & 66.4 & 66.4 & 66.4 & 66.4 \\
Split leakage & 91.4 & 79.6 & 79.6 & 79.2 & 79.0 & 79.6 \\
\bottomrule
\end{tabular}
\end{table}

The primary metrics also move with severity (Table~\ref{tab:primary-metrics}).
Split leakage is recovered exactly because the injected corruption directly
duplicates rows across split boundaries. PII saturates because the current primary
metric is the number of flagged columns rather than the PII hit rate within the
column; this is useful for safety gating but less useful for smooth severity
analysis.

\begin{table}
\centering
\caption{Mean primary metric value by fault family and severity.}
\label{tab:primary-metrics}
\scriptsize
\begin{tabular}{lrrrrrr}
\toprule
Fault family & 0 & 0.01 & 0.05 & 0.10 & 0.20 & 0.40 \\
\midrule
Drift & 0.102 & 0.099 & 0.126 & 0.208 & 0.374 & 0.646 \\
Duplicates & 0.000 & 0.010 & 0.048 & 0.091 & 0.167 & 0.286 \\
Fairness & 0.025 & 0.017 & 0.109 & 0.206 & 0.417 & 0.799 \\
Missingness & 0.000 & 0.003 & 0.015 & 0.028 & 0.054 & 0.099 \\
PII & 0.000 & 0.600 & 1.000 & 1.000 & 1.000 & 1.000 \\
Split leakage & 0.000 & 0.010 & 0.050 & 0.100 & 0.200 & 0.400 \\
\bottomrule
\end{tabular}
\end{table}

\subsection{RQ3: Recommendations, Policy Gates, and Runtime}

Recommendation coverage and policy-gate failure rates match the detection matrix
in this controlled benchmark. This alignment is expected because each injected
fault maps to a corresponding recommendation family and policy threshold. The
important governance result is that fairness disparity now participates in the
policy gate, rather than being reported only as an informational diagnostic.

The first majority-detected nonzero severity is summarized in
Table~\ref{tab:first-detected}. The result highlights two classes of faults:
critical categorical risks such as leakage and PII, which are actionable at very
low severity, and continuous degradation modes such as drift and missingness,
which are governed by warning thresholds.

\begin{table}
\centering
\caption{First majority-detected nonzero severity and clean-control detection rate.}
\label{tab:first-detected}
\begin{tabular}{lrr}
\toprule
Fault family & First majority-detected severity & Clean detection rate \\
\midrule
Drift & 0.20 & 0.00 \\
Duplicates & 0.05 & 0.00 \\
Fairness & 0.10 & 0.00 \\
Missingness & 0.20 & 0.00 \\
PII & 0.01 & 0.00 \\
Split leakage & 0.01 & 0.00 \\
\bottomrule
\end{tabular}
\end{table}

Runtime is low in this benchmark. Across 180 audits, total runtime was 41.188
seconds, mean runtime was 0.229 seconds per audit, and maximum runtime was 0.538
seconds. This suggests that the tabular analyzer is practical for small-to-medium
CI-style dataset gates. Scaling over rows, columns, modalities, and richer PII
detectors remains future work.

\section{Discussion}

\subsection{What the Results Support}

The experiment supports three immediate conclusions. First, ASTRID can detect
known injected tabular faults without false positives on clean controls under the
current benchmark configuration. Second, the health score and primary metrics
respond in the expected direction as severity increases. Third, the framework
connects measurement to action: detected faults trigger recommendations and
policy-gate failures.

These conclusions are deliberately narrower than a claim of universal dataset
reliability. The present study validates threshold behavior for controlled
tabular faults. It does not yet prove that ASTRID outperforms baseline tools, that
its thresholds are optimal, or that its score predicts downstream model
degradation across domains.

\subsection{Practical Implications}

ASTRID is most useful as an early dataset gate. Before model training, it can
identify conditions that should block or delay development: split leakage, PII,
excessive duplicates, high missingness, strong temporal drift, or subgroup
disparity. During development, it can compare dataset revisions and document
whether quality interventions improved the evidence profile. In governance
reviews, it can provide a reproducible audit trail that links raw measurements to
recommendations and policy decisions.

\subsection{Toward a Top-Tier Empirical Claim}

A stronger submission should extend this work in four directions. First, evaluate
on real public datasets where defects are injected into heterogeneous, messy data
rather than fully synthetic data. Second, compare ASTRID against category-specific
baselines such as schema validators, drift detectors, leakage checks, and
fairness audit tools. Third, connect dataset audit findings to downstream model
degradation by training fixed model families on clean and corrupted datasets.
Fourth, add time-series and image fault-injection benchmarks so the multimodal
architecture is supported by multimodal evidence.

\subsection{Limitations}
\label{sec:limits}

The current study has several limitations. The benchmark is synthetic, faults are
injected independently, and detection labels are defined using ASTRID-aligned
thresholds. The empirical scope is tabular; time-series and image reliability
claims require additional analyzers and experiments. PII is represented by
email-like strings, and the current PII primary metric saturates once a column is
flagged. Fairness is represented by positive-rate disparity, which is only one
fairness concept. Finally, the health score is a configurable decision-support
index, not a formal guarantee of safety.

\section{Conclusion}

This paper introduced ASTRID, a local-first framework for behavior-aware dataset
reliability assessment in industrial AI. ASTRID represents dataset risk as a
multi-dimensional evidence vector, computes interpretable diagnostics, and links
findings to recommendations and policy-gate decisions.

A repeated-seed tabular fault-injection benchmark shows that ASTRID detects six
controlled reliability failures at threshold-consistent severities, produces no
primary detections on clean controls, decreases health scores under injected
faults, and runs quickly enough for CI-style dataset gates on small-to-medium
tabular datasets. These results establish a credible first empirical foundation
for ASTRID and clarify the next experiments required for a full top-tier paper:
real datasets, baselines, downstream model effects, scaling studies, and
multimodal validation.

\begin{thebibliography}{12}

\bibitem{amershi2019}
Amershi, S., Begel, A., Bird, C., DeLine, R., Gall, H., Kamar, E., Nagappan, N.,
Nushi, B., Zimmermann, T.: Software engineering for machine learning: A case
study. In: ICSE-SEIP (2019)

\bibitem{breck2019}
Breck, E., Polyzotis, N., Roy, S., Whang, S., Zinkevich, M.: Data validation for
machine learning. In: SysML (2019)

\bibitem{gebru2021}
Gebru, T., Morgenstern, J., Vecchione, B., Vaughan, J.W., Wallach, H., Daume III,
H., Crawford, K.: Datasheets for datasets. Communications of the ACM 64(12),
86--92 (2021)

\bibitem{iso25012}
ISO/IEC: ISO/IEC 25012:2008 Software engineering -- Software product Quality
Requirements and Evaluation (SQuaRE) -- Data quality model. International
Organization for Standardization (2008)

\bibitem{kaufman2012}
Kaufman, S., Rosset, S., Perlich, C.: Leakage in data mining: Formulation,
detection, and avoidance. ACM Transactions on Knowledge Discovery from Data 6(4)
(2012)

\bibitem{mitchell2019}
Mitchell, M., Wu, S., Zaldivar, A., Barnes, P., Vasserman, L., Hutchinson, B.,
Spitzer, E., Raji, I.D., Gebru, T.: Model cards for model reporting. In: FAT*
(2019)

\bibitem{nist2023}
National Institute of Standards and Technology: Artificial Intelligence Risk
Management Framework (AI RMF 1.0). NIST AI 100-1 (2023)

\bibitem{polyzotis2017}
Polyzotis, N., Roy, S., Whang, S.E., Zinkevich, M.: Data management challenges
in production machine learning. In: SIGMOD (2017)

\bibitem{ribeiro2020}
Ribeiro, M.T., Wu, T., Guestrin, C., Singh, S.: Beyond accuracy: Behavioral
testing of NLP models with CheckList. In: ACL (2020)

\bibitem{sambasivan2021}
Sambasivan, N., Kapania, S., Highfill, H., Akrong, D., Paritosh, P., Aroyo, L.M.:
Everyone wants to do the model work, not the data work: Data cascades in high-
stakes AI. In: CHI (2021)

\bibitem{schelter2018}
Schelter, S., Lange, D., Schmidt, P., Celikel, M., Biessmann, F., Grafberger, A.:
Automating large-scale data quality verification. Proceedings of the VLDB
Endowment 11(12), 1781--1794 (2018)

\bibitem{sculley2015}
Sculley, D., Holt, G., Golovin, D., Davydov, E., Phillips, T., Ebner, D.,
Chaudhary, V., Young, M., Crespo, J.F., Dennison, D.: Hidden technical debt in
machine learning systems. In: NeurIPS (2015)

\end{thebibliography}

\end{document}
